\documentclass[fleqn,10pt]{wlscirep}
\usepackage[utf8]{inputenc}
\usepackage[T1]{fontenc}
\usepackage{amsmath,amssymb,amsfonts,mathtools}
\usepackage{booktabs}
\usepackage{graphicx}
\usepackage{xcolor}

\title{Semantic network structure of an AI risk taxonomy reveals loss-of-control concepts as connector hubs}

\author[1,2,*]{Youngsam Chun}

\affil[1]{Adjunct Professor, Ulsan National Institute of Science and Technology (UNIST), Ulsan, Republic of Korea}
\affil[2]{Frontier AI Lab, KT Corporation, Seoul, Republic of Korea}
\affil[*]{deep1003@snu.ac.kr; ORCID: \href{https://orcid.org/0000-0002-6877-6230}{0000-0002-6877-6230}}

\keywords{AI risk taxonomy, semantic network, complex networks, expectation--maximization, loss of control, responsible AI}

\begin{document}

\sloppy
\raggedbottom

\begin{abstract}
Operational AI risk taxonomies are trees, although risk meanings overlap
across category boundaries. We analyse the
semantic space of a large taxonomy: 1,660 active level-4 risk
cards distilled via a documented evidence chain from 3,906,767 research,
policy, and patent records. Two expectation--maximization
procedures structure the space: a constrained-EM audit stabilizing category
assignment, and a graph-regularized seeded spherical EM partitioning the
card network into 50 communities. In the released network, 27.7\% of cards
belong meaningfully to two or more level-3 families---the first
quantification of what a single-parent tree discards. To measure connectivity as a
property of the risks themselves, we rebuild the network as a
similarity-threshold graph with no nearest-neighbour floor. Degree heterogeneity is substantial (Gini
0.50), yet formal tests reject a power-law tail. The semantic
core is dominated by concepts of eroding human control and miscalibrated
trust---automation bias, overreliance, delegation to misaligned agents,
collectively inefficient multi-agent equilibria---which rank highest across
thresholds, connect over 40 of 50 communities, and are far more likely to
be under taxonomy review. This centrality is a structural property of
system-level risk concepts, with implications for taxonomy design and
review prioritization.
\end{abstract}

\maketitle

% =====================================================================
\section*{Introduction}

Organizations that build or govern AI systems increasingly rely on structured
risk taxonomies to inventory what can go wrong. Public efforts include
harm taxonomies for language models\cite{weidinger2022}, the MIT AI Risk
Repository, which consolidates hundreds of risks from dozens of source
frameworks\cite{slattery2024}, surveys of catastrophic AI
risks\cite{hendrycks2023}, and governance frameworks such as the NIST AI Risk
Management Framework\cite{nist2023} and the OECD classification of AI
systems\cite{oecd2022}. In operational practice these taxonomies share one
structural commitment: they are trees. Every leaf-level risk is filed under
exactly one parent category, because single-parent hierarchies are easy to
audit, assign ownership over, and keep consistent across releases.

Risks, however, are not defined by their position in a tree. They are defined
by natural-language descriptions of failure mechanisms, and the semantics of
those descriptions overlap freely across category boundaries. A card about
AI-generated scams is simultaneously about fraud, misinformation, and
manipulation; a card about loss of human oversight touches nearly every
operational category. A single-parent hierarchy must pick one shelf for each
such concept and silently discard the rest of its meaning. How much is
discarded, and which concepts sit where the discarding is most severe, has to
our knowledge not been quantified for a production AI risk taxonomy.

Here we treat a large operational taxonomy---the KT RAI Risk Taxonomy 2.0,
release v2.17.2---as an empirical semantic system and study it with the tools
of network science\cite{newman2018}. The taxonomy's 1,660 active level-4 (L4)
risk cards descend, through a documented and reproducible evidence chain, from
3,906,767 research papers, policy reports, and patents. Each card carries a
bilingual mechanism definition, and the released pipeline embeds these
definitions\cite{chen2024bgem3}, links each card to its nearest semantic
neighbours, and partitions the resulting proximity network into semantic
communities.

This paper makes three contributions. First, we describe the two
expectation--maximization (EM) procedures\cite{dempster1977} that give the
taxonomy its semantic structure: a \emph{constrained-EM audit} used during
construction to test and stabilize category assignments under explicit
guards, and a \emph{graph-regularized seeded spherical EM}\cite{banerjee2005}
used to detect communities in the proximity network while respecting both
embedding geometry and network topology. Second, we characterize the semantic
proximity network itself and quantify cross-category overlap: 459 cards
(27.7\%) carry non-trivial responsibility in two or more level-3 (L3)
families, which is precisely the information a single-parent tree cannot
represent. Third, we characterize the connectivity of the risk space as a
property of the risks themselves rather than of any particular graph
construction: because the released union-$k$NN network imposes a degree
floor and suppresses hubs, we recompute the card embeddings and rebuild the
network as a similarity-threshold graph with no floor, subject the resulting
degree distribution to formal tail diagnostics\cite{clauset2009,broido2019},
and show that its semantic core is dominated by concepts of eroding human
control and miscalibrated trust. We explain why such system-level concepts
must occupy the centre of any sufficiently comprehensive AI risk space,
connecting an empirical observation about taxonomy structure to theoretical
arguments about gradual human disempowerment\cite{russell2019,bengio2024,
kulveit2025}.

Throughout, we report what the tail statistics actually support: the degree
distribution is strongly heterogeneous but likelihood-ratio tests reject a
power-law tail, so we make no scale-free claim\cite{clauset2009,broido2019}.
And we treat semantic proximity as a statement about meaning, not about
causal transmission between risks\cite{helbing2013}: the network records
which failure descriptions resemble one another, not which failures cause
one another.

% =====================================================================
\section*{Results}

\subsection*{An evidence-grounded risk registry}

The analysed release (v2.17.2) rests on a five-step distillation
(Methods). An integrated AI evidence space of 3,906,767 records (1,749,159
papers, 125,372 policy reports, 2,032,236 patents) is filtered to a
source-governed AI-risk evidence view of 82,971 records. Joining source risk
identifiers and reviewing redundancy yields 1,726 candidate L4 rows;
exact-label deduplication and removal of taxonomy-level leakage yield 1,711
canonical cards with permanent identifiers; retiring merged records leaves
1,660 active cards, each carrying a mechanism-level bilingual definition and
its evidence references. All 1,660 active cards are assigned to one of the
taxonomy's L3 families, with 765 cards (46.1\%) additionally carrying a HOLD
marker indicating that the assignment is operationally usable but queued for
further human review. Card-level scalar attributes such as estimated severity
or probability exist in the registry but are intentionally excluded from all
analyses reported here, because they are editorial estimates rather than
measured quantities.

\subsection*{The semantic proximity network}

Embedding each card's bilingual mechanism definition with BGE-M3 and linking
mutual nearest neighbours (base $k=8$, minimum cosine similarity 0.62,
projected union-$k$NN $k=10$; edge weight mixes an L3-profile similarity with
direct cosine similarity at 0.65/0.35; Methods) produces a network of 1,660
nodes and 11,869 edges. The network is a single connected component with
density 0.0086, mean degree 14.30 (median 13), maximum degree 41, average
clustering coefficient 0.333, and slightly positive degree assortativity
($+0.093$). Figure~\ref{fig:map} shows the network in its released
ForceAtlas2 layout\cite{jacomy2014}. This released network defines the
taxonomy's community structure and visual layout; because its union-$k$NN
construction guarantees every node ten projected neighbours and truncates
each node's outgoing choices, we use a separate construction-free network,
introduced below, for all degree and centrality analysis.

The graph-regularized seeded spherical EM partitions the network into 50
semantic communities (sizes 2--140, median 25), each seeded by an L3-family
centroid but free to reorganize cards across family boundaries
(Methods). The largest communities are Weaponization (140 cards), Policy
Exposure (114), Non-Contestability (90), Anthropomorphism (71), and
Misinformation/Disinformation (68). The EM converged in 16 iterations to a
mean within-community cosine similarity of 0.751.

\subsection*{Quantifying what a tree discards}

For every card we compute a responsibility distribution over L3 families from
its neighbourhood structure (Methods). At a responsibility threshold of
0.03, 459 of 1,660 cards (27.7\%) belong appreciably to two or more
families, and individual cards reach up to eight family affiliations. In
other words, more than a quarter of the taxonomy's leaf concepts are
multi-parent concepts being stored in a single-parent data structure. This is,
to our knowledge, the first direct quantification of the representational
loss imposed by tree-shaped AI risk taxonomies, and it identifies exactly
which cards would benefit from an explicit multi-label extension.

Cross-community membership is the norm rather than the exception at the
level of neighbourhoods as well: in the construction-free network
introduced below, the median participation coefficient\cite{guimera2005}
across all cards is 0.880, meaning that a typical card's semantic
neighbours are spread over many communities.

\subsection*{Degree structure without construction artifacts}

Degree in the released network is not a clean property of the risks: the
union-$k$NN construction imposes a hard floor at $k=10$ and truncates every
node's outgoing choices, compressing heterogeneity (Gini 0.153) and
suppressing hubs. To measure connectivity as a characteristic of the risk
concepts themselves, we recompute the BGE-M3 embeddings of all 1,660 cards
(reproduction fidelity against the released edge similarities: $r=0.980$,
mean absolute error 0.014; Methods) and build a similarity-threshold graph:
two cards are linked if and only if their cosine similarity is at least
$\theta$, with no neighbour count imposed anywhere. At the pipeline's own
minimum-similarity value $\theta=0.62$ this yields 90,848 edges, mean degree
109.5 (median 69), maximum degree 554, five isolates, a giant component of
1,655 nodes, clustering 0.432, and assortativity $+0.197$.

Freed from the floor, the degree distribution is strongly heterogeneous:
the Gini coefficient more than triples to 0.504, the maximum degree is
eight times the median, and 162 cards (9.8\%) exceed four times the median
(Fig.~\ref{fig:structure}a). Heterogeneity is not an artifact of the
specific threshold: across $\theta \in [0.55, 0.70]$ the Gini coefficient
rises monotonically from 0.33 to 0.58 while degree rankings remain stable
(Spearman $\rho = 0.89$ between $\theta=0.62$ and $\theta=0.68$; $\rho =
0.94$ against the fully threshold-free strength $\sum_j s_{ij}$;
Fig.~\ref{fig:structure}b). Formal tail diagnostics, however, reject a
power law: likelihood-ratio tests\cite{clauset2009} favour lognormal and
exponential alternatives over a power-law tail (all $p<0.001$), consistent
with the observation that genuinely scale-free networks are
rare\cite{broido2019}. The risk space is therefore best described as
broadly---but not scale-freely---heterogeneous: a substantial semantic core
of highly connected concepts over a large periphery of specific mechanisms.

\subsection*{Concepts of eroding human control dominate the semantic core}

Table~\ref{tab:hubs} lists the ten highest-degree cards of the
construction-free network. The single most-connected concept in the entire
risk space is RAI4-0592, ``Collectively inefficient equilibria among AI
agents'' ($k=554$), and its rank is robust: it is first at every threshold
from 0.60 to 0.65 and second at 0.68. More broadly, the top of the ranking
interleaves two kinds of concepts: system-level erosion of human control
and trust---automation bias (rank 2), delegation of decision-making power
to misaligned AIs, misaligned rogue behaviour, overreliance undermining
user autonomy, anthropomorphic overtrust, miscalibrated trust, loss of
human agency and autonomy---and broad misuse and integrity concepts (abuse
and misuse, systemic bias, misinformation dissemination, insufficient
security). Coding the 30 highest-degree cards by mechanism, 16 of 30
describe the erosion of human control, autonomy, or calibrated trust over
AI systems. The loss-of-control theme is thus not carried by a single
labelled card but saturates the semantic core, and it sharpens as the
threshold tightens: at $\theta=0.68$, overreliance (RAI4-1552) ranks first
and loss of human agency (RAI4-0894) fifth.

Transparency requires noting what changed relative to the released
construction: RAI4-1616 (``Future AI systems might actively reduce human
control''), the highest-degree node of the union-$k$NN network, ranks 58th
(top 3.5\%) in the construction-free graph
(Fig.~\ref{fig:structure}d)---its former first place partly reflected
mutual-selection dynamics of the $k$NN projection rather than raw semantic
reach. The thematic conclusion survives the correction; the specific
figurehead does not, which is precisely why the construction-free analysis
matters.

These hubs are connectors, not local centres. The four highest-ranked
control-erosion cards have participation coefficients\cite{guimera2005} of
0.92--0.94 and neighbourhoods spanning 42--47 of the 50 communities;
top-decile hubs generally sit at the 87th--100th percentiles of $P$. In the
role terminology of Guimer\`a and Amaral, the semantic core of the taxonomy
is built from connector hubs: concepts whose neighbours are distributed
across many communities rather than concentrated in their own.

HOLD status rises monotonically with degree: 63.5\% of cards in the top
degree decile carry HOLD markers versus 44.1\% of the remainder, against an
overall rate of 46.1\% (Fig.~\ref{fig:structure}c). We interpret this not
as a defect of the hub cards but as a structural signal: cards whose
meaning spans many families are precisely the cards for which a single
released parent is hardest to defend, and the review process detects this.

\begin{table}[ht]
\centering
\caption{The ten highest-degree L4 risk cards in the construction-free
similarity-threshold network ($\theta=0.62$, release v2.17.2 card set).
$k$: degree; $s$: strength (summed similarity weight); community: EM
community of the card in the released network; $P$: participation
coefficient (percentile in parentheses); review: HOLD marks cards queued
for taxonomy review. Cards marked $\dagger$ describe erosion of human
control, autonomy, or calibrated trust.}
\label{tab:hubs}
\small
\begin{tabular}{llrrlrl}
\toprule
ID & Label (abbreviated) & $k$ & $s$ & Community & $P$ & review \\
\midrule
RAI4-0592$^\dagger$ & Collectively inefficient equilibria among AI agents & 554 & 366 & Conflict & 0.94 (96) & HOLD \\
RAI4-1482$^\dagger$ & Automation bias & 527 & 349 & Non-Contestability & 0.94 (97) & --- \\
RAI4-1013 & Abuse and misuse & 523 & 347 & Weaponization & 0.93 (89) & HOLD \\
RAI4-1566 & Systemic bias across specific communities & 499 & 328 & Hate and Unfairness & 0.94 (96) & --- \\
RAI4-0844 & Disseminating false or misleading information & 494 & 327 & Misinf./Disinf. & 0.94 (99) & HOLD \\
RAI4-1186 & Insufficient security measures & 489 & 319 & Weaponization & 0.93 (87) & --- \\
RAI4-1038 & Misapplication of AI & 487 & 321 & Goal Misalignment & 0.95 (99) & --- \\
RAI4-1425$^\dagger$ & Delegating decision-making power to misaligned AIs & 487 & 322 & Goal Misalignment & 0.93 (92) & --- \\
RAI4-0565$^\dagger$ & Misaligned rogue AI behavior & 485 & 318 & Goal Misalignment & 0.95 (100) & --- \\
RAI4-1719 & Lack of security and resilience & 485 & 319 & Non-Contestability & 0.93 (93) & HOLD \\
\bottomrule
\end{tabular}
\end{table}

\begin{figure}[ht]
\centering
\includegraphics[width=\linewidth]{figures/fig1_network_map.pdf}
\caption{\textbf{The semantic proximity network of 1,660 active L4 risk
cards.} Nodes are risk cards positioned by the released ForceAtlas2 layout;
node area scales with construction-free degree (similarity-threshold graph,
$\theta=0.62$); colours mark the ten largest of the 50 EM-detected semantic
communities (legend, with sizes), and grey marks the remaining 40. Edges of
the released union-$k$NN network (11,869) are drawn in light grey. Black
rings and italic labels mark four of the highest-ranked control-erosion
hubs (RAI4-0592, RAI4-1482, RAI4-1552, RAI4-0894), which sit in different
communities and regions of the space, consistent with their role as
cross-community connectors. Bold labels mark the centroids of the five
largest communities. The layout is used for display only; no analysis in
this paper depends on layout coordinates.}
\label{fig:map}
\end{figure}

\begin{figure}[ht]
\centering
\includegraphics[width=\linewidth]{figures/fig2_structure.pdf}
\caption{\textbf{Degree structure of the construction-free risk network.}
(\textbf{a}) Complementary cumulative degree distribution of the
similarity-threshold graph ($\theta=0.62$; log--log axes). The distribution
is strongly heterogeneous (median 69, maximum 554, Gini 0.50), but
likelihood-ratio tests reject a power-law tail against lognormal and
exponential alternatives (all $p<0.001$), so no scale-free claim is made.
(\textbf{b}) Threshold sensitivity: median degree, maximum degree (left
axis, log scale), and Gini coefficient (right axis) across $\theta \in
[0.55, 0.70]$. Heterogeneity grows as the threshold tightens; degree
rankings remain stable (Spearman $\rho=0.89$ between $\theta=0.62$ and
$0.68$). (\textbf{c}) Share of cards under HOLD taxonomy review by degree
decile. HOLD rises monotonically with degree, reaching 63.5\% in the top
decile against an overall rate of 46.1\%. (\textbf{d}) Released union-$k$NN
degree versus construction-free degree (Spearman $\rho=0.36$). The $k$NN
floor at 10 (dotted line) and hub suppression flatten the released degrees;
RAI4-0592 (rank 1, construction-free) and RAI4-1616, whose union-$k$NN
rank~1 falls to rank~58 without the construction, are highlighted.}
\label{fig:structure}
\end{figure}

% =====================================================================
\section*{Discussion}

\subsection*{Why loss-of-control concepts must sit at the centre}

The empirical finding is that concepts of declining human control occupy the
top of the degree ranking and act as cross-community connectors. We offer
three mutually reinforcing explanations.

First, \emph{semantic generality}. Control-erosion cards describe
system-level failure conditions---human oversight weakening, agency shifting
to machines, trust miscalibrated---rather than domain-specific mechanisms.
Almost every concrete risk description (a manipulated consumer, an
unaccountable decision, an autonomous weapon, an addictive companion)
contains, as part of its meaning, a local instance of humans losing control
over an outcome. In embedding space, the general concept therefore sits
within the similarity threshold of many specific ones, and degree records
exactly that reach. The same mechanism places broad misuse concepts high in
the ranking, which is why generality alone does not fully explain the
pattern: among comparably general concepts, those about control and trust
are systematically over-represented (16 of the top 30).

Second, \emph{bridging position}. The hubs' participation coefficients of
0.92--0.95, with neighbourhoods spanning more than 40 of the 50
communities, show that their connectivity is not concentrated in one
region: these cards lie on short semantic paths between operational
communities (manipulation, governance, safety, economic concentration). A
reader moving conceptually from ``persuasive design'' to ``concentration of
decision power'' passes naturally through ``humans no longer steer the
system''. Connector-hub status is the network signature of that conceptual
throughway.

Third, \emph{consilience with theory}. Independent theoretical work argues
that the most structurally important AI risks are not exotic single failures
but gradual erosions of human control: the assistance-game framing of
Russell\cite{russell2019}, the argument by Bengio and colleagues that
advanced agentic systems require management of catastrophic and
loss-of-control scenarios\cite{bengio2024}, and the gradual-disempowerment
analysis of Kulveit and colleagues\cite{kulveit2025}. The specific hubs we
find read like that literature's mechanism list: automation bias,
overreliance, delegation to misaligned agents, and collectively inefficient
multi-agent equilibria are precisely the incremental channels through which
gradual disempowerment is argued to proceed. Our result is an empirical,
bottom-up counterpart: when nearly seventeen hundred risk concepts
distilled from four decades of literature, policy, and patents are laid out
by meaning alone, the concepts these theories place at the centre of
concern also sit at the centre of the semantic space. The taxonomy was not
designed to produce this: community seeds, guards, and audits operate at
the family level, and no step of the pipeline privileges any card.

\subsection*{Implications for taxonomy design and review}

Two practical consequences follow. First, the 27.7\% multi-affiliation rate
argues for multi-label extensions of operational risk taxonomies: a released
single parent can remain the operational assignment while secondary family
memberships are stored alongside, exactly where our responsibility
distributions are non-trivial. Second, degree and participation supply a
principled review-prioritization signal. The monotonic rise of HOLD status
with degree---63.5\% in the top decile---indicates that the human review
process already gravitates toward semantically ambiguous, high-centrality
cards; making that signal explicit would let review effort track structural
importance directly.

\subsection*{Limitations}

Three caveats bound the interpretation. Semantic proximity is not causal
transmission: our network records similarity of failure descriptions, not
pathways along which one failure triggers another, and coupling of the
causal kind requires separate evidence of mechanism and temporal
order\cite{helbing2013}. Absolute degree values still depend on the choice
of similarity threshold; our claims therefore rest on rankings and
concentration measures shown to be stable across thresholds ($\rho \geq
0.89$) and against the fully threshold-free strength centrality ($\rho =
0.94$), not on any single graph. Degree in a similarity graph also rewards
semantic generality as such---broad catch-all labels rank high alongside
substantively central concepts---which is why we report the thematic
composition of the core rather than treating any individual rank as
meaningful by itself. Finally, the underlying registry, while
evidence-grounded and bilingual, reflects the coverage of its source
corpora; risks under-represented in research, policy, and patent literature
are under-represented here as well.

% =====================================================================
\section*{Methods}

\subsection*{Evidence chain and registry}

The integrated AI evidence space contains 3,906,767 records: 1,749,159 papers
(global Web of Science AI collections), 125,372 policy reports (curated
AI-policy master plus policy-materials collection), and 2,032,236 patents
(AI patent master). Records are standardized to a common schema, retained
only with a non-empty title, and deduplicated within document type and
analytical category by an ordered key hierarchy (patent application ID, DOI,
source record ID, URL, normalized title--year--institution fallback). A
source-governed AI-risk evidence view selects 82,971 records (50,375 papers,
938 policy reports, 31,658 patents). Joining source risk identifiers through
a crosswalk and reviewing semantic redundancy yields 1,726 candidate L4
rows; exact normalized-label deduplication (semantic-similarity thresholds
are deliberately not used) and removal of fourteen taxonomy-level leakage
rows yield 1,711 canonical cards with permanent identifiers RAI4-0001 through
RAI4-1726; registry-preserving consolidation of duplicate lineages leaves
1,660 active cards, with all retirements kept as traceable
\texttt{merged\_into} records. Card-level severity, probability, and impact
attributes are excluded from every analysis in this paper.

\subsection*{Algorithm 1: constrained-EM audit}

During taxonomy construction, category assignments are tested with a
constrained expectation--maximization audit rather than unconstrained
nearest-centroid reassignment. Each audit allows only a selected set of
source-family cards to move, treats all other cards as fixed anchors,
excludes protected destinations (the 182-card Physical AI subtree), and
alternates between (E) scoring each movable card against candidate family
centroids and (M) recomputing those centroids, until no assignment changes.
The composite score of card $i$ against family $k$ is
\begin{equation}
S(i,k) \;=\; 0.60\,c(i,k) \;+\; 0.30\,d(i,k) \;+\; 0.10\,w(i,k),
\end{equation}
where $c$ is cosine similarity to the current family centroid, $d$ is cosine
similarity to the family definition, and $w$ is TF--IDF keyword cosine
similarity. A move is accepted only if the destination improves on the source
by at least 0.020 and has keyword cosine of at least 0.015, with additional
destination-specific definition guards. Audited moves remain HOLD until
human adjudication. An unguarded control run on revised definitions produced
151 candidate moves, used as sensitivity evidence only; the guarded release
records 22 defensible proposals, of which 14 change the operational parent.

\subsection*{Semantic proximity network}

Each active card's bilingual mechanism text (English label and definition,
with Korean as supplementary signal) is embedded with BGE-M3\cite{chen2024bgem3}
(pinned snapshot) and $\ell_2$-normalized. A base mutual-$k$NN graph is built
with $k=8$ and minimum cosine similarity 0.62; a projected union-$k$NN graph
with $k=10$ defines the released edge set, so every node holds at least its
own ten projected neighbours (degree floor). Edge weights combine an
L3-profile similarity (cosine similarity between the cards'
family-responsibility profiles) and the direct embedding cosine at
$0.65/0.35$. The released two-dimensional layout is ForceAtlas2\cite{jacomy2014}
with fixed seed 42 and is used for display only.

\subsection*{Algorithm 2: graph-regularized seeded spherical EM}

Communities are detected with a spherical (von Mises--Fisher-style) EM on the
unit-normalized embeddings\cite{banerjee2005}, seeded by L3-family centroids
(seed weight 5) and regularized by the network: the E-step responsibility of
community $g$ for card $i$ combines an embedding term
$\exp\{\beta\,\mathbf{x}_i^{\top}\boldsymbol{\mu}_g\}$ ($\beta=18$) with a
graph term weighting the community memberships of $i$'s network neighbours
($\lambda=2$), and the M-step re-estimates unit-norm community centroids from
responsibility-weighted cards. The procedure converged in 16 iterations
(final change rate 0, final mean within-community cosine similarity 0.751),
yielding 50 communities of sizes 2--140. Multi-affiliation counts use the
final responsibility distributions with threshold 0.03.

\subsection*{Construction-free similarity network}

For degree and centrality analysis we rebuild the network without any
nearest-neighbour machinery. The bilingual card texts (English label and
definition, followed by the Korean label and definition) are re-embedded
with BGE-M3 (dense embeddings, FlagEmbedding implementation,
$\ell_2$-normalized) and validated against the released artefact: across
the 11,869 released edges, the recomputed cosine similarities match the
stored direct similarities with $r=0.980$ and mean absolute error 0.014.
The construction-free graph links cards $i$ and $j$ if and only if
$\cos(\mathbf{x}_i,\mathbf{x}_j) \geq \theta$, weighted by the cosine. The
primary threshold $\theta=0.62$ is the released pipeline's own
minimum-similarity parameter; sensitivity is reported for $\theta \in
\{0.55, 0.58, 0.60, 0.62, 0.65, 0.68, 0.70\}$, together with the fully
threshold-free strength centrality $s_i=\sum_j \cos(\mathbf{x}_i,
\mathbf{x}_j)$. Rank stability across these choices is summarized with
Spearman correlations.

\subsection*{Network statistics and tail diagnostics}

Statistics are computed with NetworkX: degree, strength, clustering, degree
assortativity, Gini coefficients, community size distributions, and
participation coefficients $P_i = 1-\sum_g (k_{ig}/k_i)^2$ over the 50 EM
communities\cite{guimera2005} of the released network. Percentiles for $P$
are computed over all 1,660 cards. Tail diagnostics follow Clauset and
colleagues\cite{clauset2009} as implemented in the \texttt{powerlaw}
package: a discrete power law is fitted above an estimated $x_{\min}$
(maximum-likelihood $\hat{\alpha}=1.93$, $x_{\min}=42$) and compared with
lognormal, exponential, stretched-exponential, and truncated power-law
alternatives via normalized likelihood-ratio tests; the power law is
rejected against every alternative (all $p<0.001$). Thematic coding of the
30 highest-degree cards as control-erosion versus other was performed by
the author from the card definitions and is reproducible from the released
registry. Analysis code that reproduces every number and figure in this
paper is available with the data (Data availability).

\subsection*{Use of artificial intelligence tools}

Claude (Anthropic) was used to assist with analysis scripting, figure
preparation, and manuscript editing under the author's direction. The author
reviewed and takes full responsibility for all content.

% =====================================================================
\section*{Data availability}
The complete release data for taxonomy v2.17.2, including the semantic
proximity network (\texttt{semantic\_proximity\_network.json}), card
coordinates and registry (\texttt{risk\_space.json}), and an interactive
explorer, are publicly available at
\url{https://github.com/deep1003/RAI-Risk-Taxonomy-2.0}. A versioned archive
will be deposited on Zenodo with a DOI upon acceptance.

\section*{Code availability}
The analysis scripts that reproduce all statistics and figures
(\texttt{embed\_cards.py}, \texttt{threshold\_network.py},
\texttt{network\_stats.py}, \texttt{make\_figures\_v2.py}) are available in
the same repository under the MIT licence.

\section*{Acknowledgements}
The author thanks the KT Responsible AI team for maintaining the taxonomy
release infrastructure, and the Ulsan National Institute of Science and
Technology (UNIST) for academic support.

\section*{Author contributions}
Y.C.\ designed the study, conducted all analyses, prepared the figures, and
wrote the manuscript.

\section*{Competing interests}
The author is employed by KT Corporation, which develops and maintains the
RAI Risk Taxonomy analysed in this work. The author declares no other
competing interests.

% ----- References -----
\begin{thebibliography}{20}

\bibitem{weidinger2022}
Weidinger, L. \emph{et~al.}
Taxonomy of risks posed by language models.
In \emph{Proceedings of the 2022 ACM Conference on Fairness, Accountability, and Transparency}, 214--229 (2022).

\bibitem{slattery2024}
Slattery, P. \emph{et~al.}
The AI Risk Repository: a comprehensive meta-review, database, and taxonomy of risks from artificial intelligence.
Preprint at \url{https://arxiv.org/abs/2408.12622} (2024).

\bibitem{hendrycks2023}
Hendrycks, D., Mazeika, M. \& Woodside, T.
An overview of catastrophic AI risks.
Preprint at \url{https://arxiv.org/abs/2306.12001} (2023).

\bibitem{nist2023}
National Institute of Standards and Technology.
\emph{Artificial Intelligence Risk Management Framework (AI RMF 1.0)}.
NIST AI 100-1 (2023).

\bibitem{oecd2022}
OECD.
\emph{OECD Framework for the Classification of AI Systems}.
OECD Digital Economy Papers No.~323 (2022).

\bibitem{newman2018}
Newman, M. E. J.
\emph{Networks}, 2nd edn (Oxford University Press, 2018).

\bibitem{chen2024bgem3}
Chen, J. \emph{et~al.}
BGE M3-Embedding: multi-lingual, multi-functionality, multi-granularity text embeddings through self-knowledge distillation.
Preprint at \url{https://arxiv.org/abs/2402.03216} (2024).

\bibitem{dempster1977}
Dempster, A. P., Laird, N. M. \& Rubin, D. B.
Maximum likelihood from incomplete data via the EM algorithm.
\emph{Journal of the Royal Statistical Society: Series B} \textbf{39}, 1--22 (1977).

\bibitem{banerjee2005}
Banerjee, A., Dhillon, I. S., Ghosh, J. \& Sra, S.
Clustering on the unit hypersphere using von Mises--Fisher distributions.
\emph{Journal of Machine Learning Research} \textbf{6}, 1345--1382 (2005).

\bibitem{russell2019}
Russell, S.
\emph{Human Compatible: Artificial Intelligence and the Problem of Control}
(Viking, 2019).

\bibitem{bengio2024}
Bengio, Y. \emph{et~al.}
Managing extreme AI risks amid rapid progress.
\emph{Science} \textbf{384}, 842--845 (2024).

\bibitem{kulveit2025}
Kulveit, J. \emph{et~al.}
Gradual disempowerment: systemic existential risks from incremental AI development.
Preprint at \url{https://arxiv.org/abs/2501.16946} (2025).

\bibitem{clauset2009}
Clauset, A., Shalizi, C. R. \& Newman, M. E. J.
Power-law distributions in empirical data.
\emph{SIAM Review} \textbf{51}, 661--703 (2009).

\bibitem{broido2019}
Broido, A. D. \& Clauset, A.
Scale-free networks are rare.
\emph{Nature Communications} \textbf{10}, 1017 (2019).

\bibitem{helbing2013}
Helbing, D.
Globally networked risks and how to respond.
\emph{Nature} \textbf{497}, 51--59 (2013).

\bibitem{guimera2005}
Guimer\`a, R. \& Amaral, L. A. N.
Functional cartography of complex metabolic networks.
\emph{Nature} \textbf{433}, 895--900 (2005).

\bibitem{jacomy2014}
Jacomy, M., Venturini, T., Heymann, S. \& Bastian, M.
ForceAtlas2, a continuous graph layout algorithm for handy network visualization designed for the Gephi software.
\emph{PLOS ONE} \textbf{9}, e98679 (2014).

\end{thebibliography}

\end{document}
